\section{Introduction}\label{s1_intro:sec}
% Main text, Section 1 (Introduction).
% Paths in "% src:" comments are relative to the repository root.

Recorded clusters of SARS-CoV-2 infection arose overwhelmingly indoors
\citep{qian2021indoor,leclerc2020settings}, and in several well-documented events the cases were not spread
evenly over the room. They sat on one side of a call-centre floor \citep{park2020}, close to the index case in
an aircraft cabin \citep{khanh2020} or on high-speed trains \citep{hu2021}, or at neighbouring restaurant
tables in the air stream of one air conditioner \citep{lu2020,li2021}.
The standard quantitative tool for indoor risk is the Wells--Riley equation
\citep{wells1955airborne,riley1978,szeto2010review}. It and the models built on it
\citep{noakes2006modelling,bazant2021guideline} treat the room air as well mixed, so that every susceptible
occupant carries the same risk.
This article examines what a minimal spatial model adds to that picture: a
susceptible--infective--recovered (SIR) epidemic among people who perform independent random walks on a
lattice of floor cells, with a mobility $D_x$ and an infection efficiency $q_x$ that depend on the cell
(workstations, corridors, meeting rooms, barriers).

\paragraph{Prior work.}
Epidemics among moving individuals have been modelled for more than three decades: lattice automata whose
individuals move between sites \citep{boccara1992automata,rhodes1996persistence}, agents moving in continuous
space \citep{gonzalez2004scaling,frasca2006dynamical,buscarino2008disease}, and reaction--diffusion processes
on metapopulation networks \citep{colizza2007reaction,belik2011natural,wang2014spatial}.
Closest to the present model is the random-walk theory of transmission between two walkers, in which infection
is a reaction that takes place at a finite rate when an infected and a susceptible walker meet
\citep{kenkre2014theory,sugaya2018analysis,kenkre2021theory}. \citet{giuggioli2022spatio} formulate it on a
lattice with exact propagators \citep{giuggioli2020exact,sarvaharman2023particle} and predict a basic
reproduction number analytically for moving susceptible, infected and recovered individuals.
The many-walker system itself is a classical object of probability and of statistical physics. For
independent continuous-time random walkers on $\mathbb{Z}^d$, with infection when an infected and a healthy
walker meet, \citet{kesten2005spread,kesten2008shape} proved that the infected set grows linearly in time with
an asymptotic shape, and, when infected walkers recover and can be infected again, that there is a critical
recovery rate below which the infection survives with positive probability and above which it dies out
\citep{kesten2006phase}. The same process is studied in physics as the diffusive epidemic process, whose
absorbing-state transition has been analysed by mean-field theory and simulation \citep{maia2007diffusive};
with permanent immunity, as here, simulations on the square lattice show a transition of dynamic-percolation
type controlled by the density of walkers \citep{carvalho2025generalized}. These results concern infinite or
large homogeneous lattices. They establish that a threshold exists for walkers; they do not give its value
for the small heterogeneous rooms considered here.
% src: data/refs_search.json (kesten2005spread, kesten2006phase, kesten2008shape, maia2007diffusive, carvalho2025generalized: Crossref records, abstracts in manual_checks; search queries of 1 October 2026)

For compartmental models the basic reproduction number $\Rz$ is the spectral radius of the next-generation
operator \citep{diekmann1990definition}, $\Rz=\srad(\Fm\Vm^{-1})$, and it is a threshold
\citep{vandendriessche2002reproduction}. Linearised at the disease-free state, the mean-field version of the
lattice model is an epidemic patch model with standard incidence, and the relevant theorems are published.
For the patch model with symmetric movement, \citet{allen2007asymptotic} defined $\Rz=\srad(\Fm\Vm^{-1})$,
proved the threshold property, showed that the growth bound $\sabs(\Fm-\Vm)$ decreases with the dispersal rate
between its two limits, and conjectured the same monotonicity for $\Rz$; the conjecture was proved by
\citet{gao2019travel} for symmetric movement and by \citet{gao2020fast} in general (see also
\citet{chen2020asymptotic}), and it is an instance of the reduction principle of
\citet{altenberg2012resolvent}. That $\Rz$ lies between the smallest and the largest patch reproduction
number, and does not depend on movement when they coincide, is due to \citet{gao2011sis}; the sharper lower
bound, which is the fast-dispersal limit, is due to \citet{gao2020fast}, and the approach to that limit at
order one over the dispersal rate was obtained by \citet{tien2015disease}, for a network on which the pathogen
rather than the hosts moves, with a next-generation matrix of the same form. Continuum counterparts are given
by \citet{allen2008asymptotic} and \citet{wang2012basic}. (We read the published versions of
\citet{gao2011sis} and \citet{tien2015disease}, whose Proposition~2.2 and Section~4.2 contain the results
attributed to them, and of \citet{gao2020fast}. For \citet{allen2007asymptotic} and \citet{gao2019travel} we
read only the abstracts; what we attribute to them follows the account of \citet{gao2020fast} and agrees with
the abstracts as far as they go: the threshold property in the first, the monotone decrease in the dispersal
rate for symmetric movement in the second.)
% src: data/refs_search.json (gao2020fast_published: Theorems 2.1, 2.3, 3.3, Corollary 3.5 "generalization of Lemma 3.4 in Allen et al.", Discussion on Tien et al.; checked 1 October 2026)
% src: published versions of doi:10.1016/j.mbs.2011.05.001 (Proposition 2.2: max{max_i tilde R0^(i), min_i R0^(i)} <= R0 <= max_i R0^(i)) and doi:10.1007/s00285-014-0791-x (Section 4.2, eqs. (29)-(35): first-order Laurent correction of R0 in 1/d_W with the generalised group inverse X_0), author copies read 2 October 2026
% src: abstracts of doi:10.1137/060672522 and doi:10.1137/18M1211957 in their Crossref records (read 2 October 2026)

Spatially resolved indoor models also exist: zonal ventilation models \citep{noakes2009mathematical},
transport models for the airborne concentration \citep{lau2022predicting,vuorinen2020modelling}, and
pedestrian- or agent-based simulations of supermarkets and of air travel
\citep{ying2021modelling,namilae2017multiscale}. They resolve the air or the trajectories in more detail than a
lattice walk does. What the lattice model offers in return is analytical access to reproduction numbers; what
it lacks, as Section~\ref{s8_outbreaks:sec} shows, is an adequate description of seated people who share air
and of who meets whom.

Three bodies of measurement bear on such a model, and each has its own literature. Close-range contacts
between people indoors have been recorded with wearable sensors in schools, offices, hospitals and
conferences; they are concentrated on few, repeated partners, structured by groups and broadly distributed in
duration \citep{cattuto2010dynamics,stehle2011high,genois2015data,genois2018can}; random walkers in a plane
that slow down near attractive individuals reproduce several of these statistics
\citep{starnini2013modeling}; and epidemics have been simulated on the recorded networks
\citep{stehle2011simulation,gemmetto2014mitigation}. The transport of
aerosol in vehicles has been measured with tracers \citep{kinahan2021aerosol,woodward2022evaluation,ou2022},
and turbulent diffusion coefficients of rooms have been related to their ventilation
\citep{cheng2011modeling}. And outbreak investigations report risk by distance: contact tracing on aircraft
uses a zone of two rows around a case, a rule whose limits are known \citep{hertzberg2016risk}, and
transmission of influenza in schools has been resolved into classes, grades and schools
\citep{endo2021within}. These are the data against which the model is tested in
Section~\ref{s8_outbreaks:sec}; the rules they suggest, a constant ratio of near to far risk and a fixed share
of contact within the group, are the generic competitors that the model has to beat.

\paragraph{Contributions.}
\begin{enumerate}[label=(\roman*),leftmargin=2.2em,itemsep=1pt,topsep=3pt]
\item \emph{Mean-field reproduction number} (Sections~\ref{s3_r0:sec} and~\ref{s4_geometry:sec}). We state
  and prove that $\Rz=\srad\bigl(\beta\Qm(\gamma\Id-\gen)^{-1}\bigr)$, a next-generation matrix built from the
  Laplace-transformed lattice propagator. For same-cell transmission we prove that
  $\beta\qmean/\gamma\le\Rz\le\beta\qmax/\gamma$, that $\Rz$ decreases monotonically with the mobility scale
  between these limits, and that any high-risk zone contributes a lower bound; the lower bound and the
  monotonicity can fail for other contact kernels. In a closed room with uniform $q$, $\Rz$ does not depend on
  the size or shape of the room. The genuine boundary effect is a leak (doors, finite dwell
  time): for uniform $q$ and the same-cell kernel, or for a dwell time that does not depend on position and
  any kernel, it replaces $\gamma$ by $\gamma+\muk$, and for heterogeneous $q$ and the same-cell kernel it
  gives $\beta\langle q\rangle_\nu/(\gamma+\muk)\le\Rroom\le\beta\qmax/(\gamma+\muk)$
  (Theorem~\ref{s4_geometry:thm:leaky}); for other kernels and a localised door the formula and both bounds
  can fail.
\item \emph{Discrete individuals} (Section~\ref{s5_pair:sec}). Following the two-walker theory of
  transmission \citep{kenkre2014theory,giuggioli2022spatio}, infection is treated as a partially absorbing
  defect of the two-walker process. This gives the expected number of secondary cases $\Rone(x)$ of an index
  case and a closed form for homogeneous periodic rooms and the infinite floor, which approximates rooms with
  reflecting walls to within 1\,\% in the cases tested. It explains why counted secondary cases fall below
  $\Rz$ at low occupancy and what form a finite-size effect takes for discrete people.
  % src: data/s5_pair_checks.json (finite_size: largest deviation of the closed form from the exact pair solve 0.90 %, 21 rooms)
\item \emph{Exact stochastic simulation} (Section~\ref{s6_scenes:sec}) of four scenes (office, supermarket,
  classroom, metro carriage) and of interventions, in which secondary cases are counted from the infection
  tree and not fitted.
\item \emph{Confrontation with data} (Section~\ref{s8_outbreaks:sec}). We assemble 38
  records of documented indoor outbreaks and exposure cohorts: a first dataset of 27 records from 31
  published sources (23 with counts of cases and of exposed people) and a second dataset of eleven records of ten
  outbreaks with seat or position data. Two spatial models ($\Mone$, $\Mtwo$) are tested against a well-mixed
  model ($\Mzero$) on 18 of them in two out-of-sample tests (two records for calibration and five held out;
  then six and five; four further records enter the level checks and negative controls of the first test) and, in the second test by design, against a constant ratio of near to far risk. Three
  further analyses address what outbreak counts cannot: contacts generated by the lattice walk are compared
  with recorded contacts in six indoor settings; the airborne kernel is fixed by published tracer measurements
  and then confronted with the outbreaks; and the model reduced to zones, with mixing measured in one school,
  is used to predict an influenza epidemic in 29 others. Each of the five primary tests follows a protocol
  written before the model under test was run on its data (for the tracer analysis, on outbreaks whose outcome
  under the other models was already known). The protocols are pre-specified, not pre-registered in the usual
  sense: they were written with the data already known, they were not lodged with a registry, and their timing
  is attested by file hashes and file times only. Comparisons added after the results were known are marked
  as post hoc, and a register of all tests says which were specified in advance
  (Supplementary Table~\ref{S-appS_prot2:tab:register}). Everything is reported with its failures.
  % src: data/bsc_outbreaks/outbreaks.csv (27 rows); code/bsc_validation2/a1_more_outbreaks/PREREGISTRATION.md section 2.1; notes/VERIFICATION.md (Section 4.4, outbreak dataset); notes/VERIFICATION.md (Section 4.5, first outbreak test: pre-registration data-aware, model-blind); notes/VERIFICATION.md section 4.10 (reservations)
\end{enumerate}

\paragraph{What is known and what is added.}
Table~\ref{s1_intro:tab:known} separates the two. The theorems on the mean-field $\Rz$ are known for patch
models. We give complete proofs from standard facts on non-negative matrices (Supplementary
Section~\ref{S-appS_r0:sec}) because the conditions matter: the lower bound $\beta\qmean/\gamma\le\Rz$ holds for
the same-cell kernel and fails for general contact kernels, for which we give a three-cell counter-example
(Proposition~\ref{s3_r0:prop:kernel}).
% src: data/plan_theory_checks.json (checks: kernel_counterexample, 0.50 against mean q 0.667)
The two-walker formalism is prior work, and the saturation of contacts between neighbours that it captures
is a classical effect \citep{ball1997epidemics,keeling1999effects,durrett1994importance}. In particular,
\citet{giuggioli2022spatio} already derive the mean number of secondary infections of one infected walker
among uniformly placed susceptible walkers in a closed domain from the two-walker formalism and check it
against simulation; the one-defect solution is that of \citet{szabo1984localized}. What is added here is
narrower (Section~\ref{s5_pair:sec}): the normalisation of the pair hazard by the occupancy, which makes the
pair-level and the mean-field numbers refer to the same model and gives the entrywise comparison
$\ngmone\le\ngm$; the numbers $\Rone(x)$ and the matrix $\ngmone$ for heterogeneous rooms with contact
kernels; the closed form for periodic rooms and the infinite floor with its limits; the comparison with the
second generation; and the observation that neither $\Ronebar$ nor the spectral radius of $\ngmone$ is a
threshold quantity (Remark~\ref{s5_pair:rem:threshold}). These statements of what is added rest on a targeted
search (Crossref and arXiv queries on infection among random walkers, the diffusive epidemic process and
patch-model reproduction numbers, 1~October 2026), not on a systematic review.
% src: data/refs_search.json (searches: queries and numbers of records screened)

\begin{table}[tbp]
\centering\small
\caption{Results of this article, where they are known from, and what is added here. `Thm' numbers refer
to the published version of \citet{gao2020fast}. No novelty is claimed for rows 1--6 and~8; in rows 9
and~10 the data are those of the cited sources; what is added is their compilation and digitisation and the
tests.}
\label{s1_intro:tab:known}
\begin{tabular}{@{}>{\raggedright\arraybackslash}p{0.27\textwidth}>{\raggedright\arraybackslash}p{0.31\textwidth}>{\raggedright\arraybackslash}p{0.35\textwidth}@{}}
\toprule
Result & Known from & Added here \\
\midrule
1.\ $\Rz=\srad(\Fm\Vm^{-1})$ and its threshold property &
  \citep{diekmann1990definition,vandendriessche2002reproduction}; for the patch model
  \citep{allen2007asymptotic} &
  proof in lattice notation; the modal formula, of which the largest diagonal term is a lower bound \\\addlinespace[2pt]
2.\ $\beta\qmean/\gamma\le\Rz\le\beta\qmax/\gamma$ &
  between the smallest and largest patch numbers \citep{gao2011sis}; lower bound \citep{gao2020fast} Thm~3.3;
  continuum case \citep{allen2008asymptotic}; heterogeneity raises $\Rz$ \citep{hasibeder1988population} &
  the condition (same-cell kernel) and a counter-example for other contact kernels \\\addlinespace[2pt]
3.\ $\Rz$ decreases with mobility, between the limits of row~2 &
  conjectured in \citep{allen2007asymptotic}; proved in \citep{gao2019travel} (symmetric movement) and
  \citep{gao2020fast} Thm~2.3 (general); \citep{chen2020asymptotic,altenberg2012resolvent} &
  size of the effect in four indoor scenes, at low and at realistic mobility \\\addlinespace[2pt]
4.\ Uniform $q$: $\Rz=\beta q/\gamma$ whatever the movement &
  \citep{gao2011sis}; loss through absorbing boundaries \citep{skellam1951random,cantrell2003spatial} &
  consequence for rooms (no size effect when closed); leaky-room formula $\Rroom=\beta q/(\gamma+\muk)$
  (same-cell kernel) \\\addlinespace[2pt]
5.\ Zone lower bound; expansions for fast and slow mixing; edge-wise monotonicity &
  fast-dispersal expansion \citep{tien2015disease}; the others are elementary consequences of the same matrix
  theory, for which a targeted search found no source &
  statements and proofs for indoor lattices \\\addlinespace[2pt]
6.\ Hotspot maps (Perron vectors, elasticities) &
  standard eigenvector sensitivity &
  application to the scenes; a caveat on transients \\\addlinespace[2pt]
7.\ Infection as a reaction between two walkers; mean number of secondary infections of one infected
  walker &
  \citep{kenkre2014theory,sugaya2018analysis,giuggioli2022spatio,kenkre2021theory}; one defect
  \citep{szabo1984localized}; local saturation
  \citep{ball1997epidemics,keeling1999effects,durrett1994importance}; thresholds for walkers on
  infinite lattices \citep{kesten2006phase,maia2007diffusive,carvalho2025generalized} &
  occupancy-normalised hazard and $\ngmone\le\ngm$; $\Rone(x)$ and $\ngmone$ for heterogeneous rooms with
  kernels; closed form for periodic rooms and the infinite floor; second generation; not a threshold
  \\\addlinespace[2pt]
8.\ Diffusing-aerosol exposure term; decline of risk with distance &
  \citep{lau2022predicting,venkatram2021modeling,cheng2011modeling}; two-row rule on aircraft
  \citep{hertzberg2016risk} &
  used only as the extension tested in Section~\ref{s8_outbreaks:sec}; no claim that it predicts better than a
  constant ratio \\\addlinespace[2pt]
9.\ Outbreak investigations &
  the published sources of Supplementary Section~\ref{S-appB_dataset:sec}; an earlier compilation of outbreaks
  evaluated with a well-mixed model \citep{peng2022practical} &
  38 records: near/far strata for eight records of the first dataset, the seat map of the call centre, the
  seat-offset matrix of the train cohort, and eleven records of the second dataset with seat or position data;
  two pre-specified out-of-sample tests and their outcome \\\addlinespace[2pt]
10.\ Recorded contacts, tracer measurements, class structure of school epidemics &
  \citep{cattuto2010dynamics,stehle2011high,genois2018can,kinahan2021aerosol,woodward2022evaluation,endo2021within};
  walkers that slow down near others \citep{starnini2013modeling} &
  pre-specified tests of the walk's contacts, of the kernel with a measured coefficient and of the zone-level
  reduction against these data, with negative or weak outcomes \\
\bottomrule
\end{tabular}
\end{table}

\paragraph{Findings.}
In the office scene (234 walkable cells, $\beta=0.5\perday$, $\gamma=0.14\perday$, mobility scale
$\Dz=1\cellday$) the mean-field number is $\Rz=5.11$, inside the bounds 4.64--5.36 that hold for every layout
with the same zone fractions.
% src: data/plan_theory_checks.json (office: D0=1.R0 = 5.107; lower_bound 4.643, upper_bound 5.357)
With 100 discrete people, however, a single index case infects on average $\Ronebar=2.27$ others (the exact
pair-level expectation; counted $\Rhat=2.265\pm0.004$, mean $\pm$ standard error over 400{,}000 runs).
% src: data/s5_pair_large_sample.json (office|D0=1: R1bar_exact 2.2661; mean 2.2648, se 0.0039, n_runs 400000)
Half of the introductions die out: the probability of an outbreak reaching 10\,\% of the room is 0.50,
against 0.78 for the mean-field branching process.
% src: data/bsc_sim/03_scenes.csv (office|D0=1|range1m: p_major 0.496 [0.474, 0.518], p_major_branching 0.781)
The two numbers respond to mobility in opposite ways: at $\Dz=10$ and $100\cellday$, $\Rz$ falls to 4.68 and
4.65 while the pair-level number rises to $\Ronebar=3.79$ and $4.34$.
% src: data/s5_pair_large_sample.json (office|D0=10: R0_ngm 4.676, R1bar_exact 3.7934; office|D0=100: R0_ngm 4.646, R1bar_exact 4.3411)
Barriers, separation of zones and capacity limits reduce simulated outbreaks at low mobility and occupancy.
The cause is the discreteness of individuals (an infective can infect each neighbour only once and keeps
meeting the same ones), not a lower $\Rz$: under frequency-dependent transmission a capacity limit leaves
$\Rz$ unchanged, separation of zones raises it, barrier cells leave it unchanged in a room with uniform $q$
(Theorem~\ref{s4_geometry:thm:closed}), and walls between cells cannot lower it
(Theorem~\ref{s3_r0:thm:edges}).
% src: data/s7_interventions_barriers.json (D=1 paired rhoB=0.10|FD attack -0.040 +- 0.003); data/bsc_sim/05_heterogeneity.json (attack_all 0.347 -> 0.231, R0_ngm 4.286 -> 4.972); data/bsc_sim/07_occupancy.json
These effects belong to a narrow physical regime. With 1.5\,m cells, $\Dz=1\cellday$ is a root-mean-square
displacement of about 3\,m in a day; people who walk 5\,\% of the time or all of the time correspond to
$2.4\times10^{3}$ to $4.9\times10^{4}\cellday$, at which the $\Rz$ of the closed office exceeds its well-mixed
value $\beta\langle q\rangle/\gamma=4.64$ by less than 0.003\,\%.
% src: data/plan_theory_checks.json (office.mobility_table: D0 = 2400: +0.0027 %; D0 = 48700: +0.0001 %); data/bsc_theory/worked_examples.json (E9_physical_D: 2433.6 and 48672 cell^2/day)
Residence time matters as well: if the office is occupied 8 hours a day instead of 24, an index case infects
on average 0.89 to 0.93 people at $\Dz=1$ and 1.34 to 1.40 at $\Dz=10\cellday$.
% src: data/s7_interventions_schedule.json (duty_exact|office|D0=1: exact_uniform_phase 0.8892, exact_start 0.9296; duty_exact|office|D0=10: 1.3359, 1.3975)

Confronted with data, the model does not show general agreement, and what it gets right is not specific to
it (Section~\ref{s8_outbreaks:sec}). The formulation with transmission only between people on the same
lattice site is not supported. An extension with a diffusing-aerosol term predicts the decline of risk with
distance better than a well-mixed model, but not better than a constant ratio of near to far risk; it fails
on several events and admits no common mixing coefficient, and the coefficients fitted to the first dataset lie
far below those measured with tracers in vehicles. Mixing shares measured in one school predict an influenza epidemic in 29 others only as well as a generic model with
one fitted ratio. The data support a smooth decline of risk with distance and a zone structure dominated by
the group, statements that the model shares with simpler ones; attack-rate levels are fitted, not predicted;
and because people are seated or at fixed stations in every outbreak that publishes positions, the spatial
heterogeneity of mobility remains untested. We present the datasets and the tests as contributions and make
no claim of agreement.
% src: notes/VERIFICATION.md (Section 4.5, first outbreak test, summary); notes/VERIFICATION.md (Section 4.10); notes/VERIFICATION.md section 4.10 (statement)

\paragraph{How the results were checked.}
Each strand of this work (the theory, the simulations, the outbreak dataset, the two outbreak tests and the
analyses of contact records, tracer measurements and zone structure) was re-checked with separately written
code, within the same project and on the same machine: the second pass re-derived the proofs, re-implemented
the computations, including a second exact simulator, retyped the outbreak counts from the sources and
searched for counter-examples. It was not carried out by anyone independent of the project, and it is not an
external review; where it re-ran analyses, compared kernels or added comparisons, it reused data loaders, event
geometries, the transcribed seat maps and, for timing benchmarks and schedules, the simulator of the analyses
(Supplementary Section~\ref{S-appA_numerics:sec:scope}). Below, `the re-check' refers to that second pass and `second simulator' to its simulator.
Where the re-check corrected a result, the corrected result is the one reported; results that it did not
recompute, and comparisons that it added after the results were known, are identified as such. What each
re-check covered is described in Supplementary Section~\ref{S-appA_numerics:sec}.
% src: notes/VERIFICATION.md (Section 4.1, theory, Section 4.2, simulations); notes/VERIFICATION.md (Section 4.4, outbreak dataset); notes/VERIFICATION.md (Section 4.5, first outbreak test); notes/VERIFICATION.md (Section 4.6, second outbreak test, Section 4.7, contact records, Section 4.8, tracer measurements, Section 4.9, zone level)

\paragraph{Outline.}
Section~\ref{s2_model:sec} defines the model, the four scenes and the parameter regime.
Section~\ref{s3_r0:sec} treats the basic reproduction number of the mean-field model and
Section~\ref{s4_geometry:sec} room size, leaks and hotspot maps. Section~\ref{s5_pair:sec} develops the
pair-level theory for discrete individuals, and Section~\ref{s6_scenes:sec} reports the simulations of the
scenes and of interventions. Section~\ref{s8_outbreaks:sec} confronts the model with documented outbreaks,
recorded contacts, tracer measurements and a school epidemic, and Section~\ref{s9_discussion:sec} discusses
what stands and the limits. The Supplementary Material holds the proofs, the numerical verification, the
outbreak datasets, the protocols, extended tables and the register of all tests; its sections, equations,
figures and tables are numbered with the prefix S.
